\documentclass[runningheads]{llncs}

\usepackage{graphicx}
\usepackage{booktabs}
\usepackage{url}
\usepackage[hidelinks]{hyperref}

\title{Monte Carlo Control vs Sarsa($\lambda$) on Text Flappy Bird}

\author{Supriya Malla}

\institute{3MD3220 -- Reinforcement Learning}

\begin{document}

\maketitle

\begin{abstract}
This report compares two tabular reinforcement learning agents on the simple
Text Flappy Bird environment: first-visit Monte Carlo control and
Sarsa($\lambda$). I use the distance-based state because it is small enough for
a value-table approach. In my experiments, both agents learn, Monte Carlo
improves faster early in training, and Sarsa($\lambda$) reaches the best final
score with a longer training budget.
\keywords{reinforcement learning \and Monte Carlo control \and Sarsa($\lambda$)}
\end{abstract}

\section{Setup}

I used the simple \texttt{TextFlappyBird-v0} environment from the course
repository~\cite{textflappybird}. The state is $(dx,dy)$, where $dx$ is the
horizontal distance to the next pipe gap and $dy$ is the vertical distance
between the bird and the center of that gap. The action space has two actions:
idle and flap.

All main experiments use the default configuration
\texttt{height=15, width=20, pipe\_gap=4}. Both agents were trained for 2000
episodes, with a maximum of 200 steps per episode, and greedy evaluation every
200 episodes on 30 episodes. I measured return and score.

The Monte Carlo agent used $\gamma=0.99$, $\epsilon=0.20$,
$\epsilon$ decay $=0.999$, and minimum $\epsilon=0.01$.
The Sarsa($\lambda$) agent used $\alpha=0.15$, $\lambda=0.9$,
$\gamma=0.99$, $\epsilon=0.10$, $\epsilon$ decay $=0.999$, and minimum
$\epsilon=0.02$.

The submitted notebook is available on GitHub.\footnote{\url{https://github.com/supu02/supriya-malla-text-flappy-bird-rl}}

\section{Results}

Table~\ref{tab:final-results} shows the final evaluation after 2000 episodes.
Both agents learned useful policies. Sarsa($\lambda$) reached the best final
result, with mean return 175.07 and mean score 16.50. Monte Carlo was close
behind, with mean return 168.70 and score 16.00.

\begin{table}[t]
\caption{Final greedy evaluation after 2000 episodes.}
\label{tab:final-results}
\centering
\small
\begin{tabular}{lcc}
\toprule
Agent & Mean Return & Mean Score \\
\midrule
Monte Carlo control & 168.70 & 16.00 \\
Sarsa($\lambda$) & 175.07 & 16.50 \\
\bottomrule
\end{tabular}
\end{table}

The learning curves in Fig.~\ref{fig:training-compare} show the main
difference. Monte Carlo starts improving earlier and stays ahead during the
first part of training. Sarsa($\lambda$) improves more slowly at first, but by
the end of training it slightly overtakes Monte Carlo in both return and score.

\begin{figure}[t]
\centering
\includegraphics[width=\textwidth]{results/report_training_compare.png}
\caption{Comparison of the training curves for both agents. The lines are
100-episode moving averages and the markers are greedy evaluations.}
\label{fig:training-compare}
\end{figure}

I also looked at the learned action-values and a small parameter sweep. Figure
~\ref{fig:q-compare} shows the learned $Q(s,\mathrm{flap})$ maps for both
agents together with the sweep results. In both agents, the high-value states
form a band around the safe region near the next gap.

\begin{figure}[t]
\centering
\includegraphics[width=\textwidth]{results/report_q_and_sweeps.png}
\caption{Action-value summaries and compact parameter sweeps. Top: learned
$Q(s,\mathrm{flap})$ maps. Bottom: score from the tested sweep settings.}
\label{fig:q-compare}
\end{figure}

\section{Discussion}

The sweep results were useful. For Monte Carlo, higher initial exploration
helped: the best tested setting was $\epsilon=0.20$, with mean return 175.40
and score 16.55. For Sarsa($\lambda$), the parameter choice mattered more. A
weak setting with $\alpha=0.05$ and $\lambda=0.7$ almost did not learn, while
$\alpha=0.15$ and $\lambda=0.9$ worked much better and reached mean return
130.8 with score 12.0.

The two versions of the Text Flappy Bird environment are useful for different
reasons. The simple version is suitable for tabular RL because the state is
small. The screen version keeps more information, but the state space is much
larger, so it would be a better fit for function approximation or deep RL.
The same issue appears in the original Flappy Bird gym~\cite{flappybirdgym}:
the RGB version is too large for a value table, and even the simple original
environment would need discretization because its observations are continuous.

I also tested transfer to a few modified level settings. Monte Carlo scored
14.77 on the training layout, then dropped to 9.67 on a taller map and 5.93 on
a wider map, but increased to 17.13 when the pipe gap was larger. Sarsa showed
similar behavior, with 17.73 on the base layout, 9.50 on the taller map, 9.57
on the wider map, and 17.50 on the larger-gap setting. So both agents
generalize partly, but both are still adapted to the training geometry.

\section{Conclusion}

Overall, both methods learned a reasonable policy on the simple Text Flappy
Bird environment. In my experiments, Monte Carlo was stronger early in
training, but Sarsa($\lambda$) became slightly better with the longer 2000
episode budget. Sarsa($\lambda$) was still more sensitive to tuning. The
simple $(dx,dy)$ environment works for tabular RL, while richer versions of
the game would likely need different methods.

\begin{thebibliography}{8}

\bibitem{textflappybird}
Christodoulidis, S.:
Text Flappy Bird Gym.
\url{https://gitlab-research.centralesupelec.fr/stergios.christodoulidis/text-flappy-bird-gym}

\bibitem{suttonbarto}
Sutton, R.S., Barto, A.G.:
Reinforcement Learning: An Introduction.
MIT Press, 2nd edn. (2018)

\bibitem{flappybirdgym}
Nogueira, G.:
Flappy Bird Gym.
\url{https://github.com/Talendar/flappy-bird-gym}

\end{thebibliography}

\end{document}
